\section{Introduction}
\label{sec:intro}

The emergence of diffusion models \cite{DDPM_paper,DDIM_paper,stable_diff,Dalle2_paper} has sparked significant interest in text-guided image synthesis and manipulation.
%In recent years, with the raise of Diffusion Models \cite{DDPM_paper,DDIM_paper,stable_diff,Dalle2_paper}, the task of text-guided image synthesis and manipulation gained enormous attention from the community. 
Building on the success in image generation, they have been extended to text-guided video generation  \cite{video-diffusion-models,Modelscope,AlignYourLatents,chen2023seine,make-a-video,girdhar2023emu,SVD,guo2023animatediff,li2023videogen,zhang2023show1,guo2023sparsectrl,text2video-zero,villegas2022phenaki}.
%The huge success in image generation led to the further extension of diffusion models to generate videos conditioned by textual prompts \cite{video-diffusion-models,Modelscope,AlignYourLatents,chen2023seine,make-a-video,girdhar2023emu,SVD,guo2023animatediff,li2023videogen,zhang2023show1,guo2023sparsectrl,text2video-zero,villegas2022phenaki}.

Despite the impressive generation quality and text alignment, the majority of existing approaches such as  \cite{video-diffusion-models,Modelscope,AlignYourLatents,SVD,zhang2023show1,opensora,opensoraplan} are mostly focused on generating short frame sequences (typically of $16$, $24$, or recently $384$ frame-length).  
%Recent developments utilize temporal video compressors to extend video generation to up to 384 frames (\cite{opensora,opensoraplan}), but exhibit  limited motion.
However, short videos generators are limited in real-world use-cases such as ad making, storytelling, etc. 

The na\"ive approach of training video generators on long videos (\eg $\geq 1200$ frames) is usually impractical. Even for generating short sequences, it typically  requires expensive training (\eg $260K$ steps and $4.5K$ batch size  in order to generate $16$ frames \cite{Modelscope}). 
%When trained only on short videos, video quality commonly degrades  when generators are made to output long videos (see Sect. \cite{?}).

Some approaches  \cite{oh2023mtvg,AlignYourLatents,video-diffusion-models,opensora}
thus extend baselines by autoregressively generating short videos based on the last frame(s) of the preceding chunk.
%
Yet, simply concatenating the noisy latents of a video chunk with the last frame(s) of the preceding chunk leads to poor conditioning with inconsistent scene transitions 
%% EXTERNAL REFERENCE ATTENTION
% (see \secrefcustom{sec:appendix.ablation_studies}). 
(see Sec.\ \textcolor{linkblue}{A.3}).
Some works \cite{SVD,2023i2vgenxl,wang2024videocomposer,dai2023animateanything,xing2023dynamicrafter} integrate also CLIP \cite{CLIP_paper} image embeddings of the last frame of the preceding chunk, which slightly improves consistency. However, they are still prone to inconsistencies across chunks 
%% EXTERNAL REFERENCE ATTENTION
% (see \figrefcustom{fig:dynamicrafter-inconsistencies})
(see Fig.\ \textcolor{linkblue}{A.7})
due to the CLIP image encoder losing crucial information necessary for accurately reconstructing the conditional frames.
The concurrent work SparseCtrl \cite{guo2023sparsectrl} utilizes a sparse encoder for conditioning. 
To match the size of the inputs, its architecture requires to concatenate additional zero-filled frames to the conditioning frames before being plugged into sparse encoder. However, this inconsistency in the input leads to inconsistencies in the output (see \secrefcustom{sec:exp_comparison_with_baselines}). 

Our experiments (see \secrefcustom{sec:exp_comparison_with_baselines}) reveal that in fact all assessed image-to-video methods produce video stagnation or strong quality degradation when applied autoregressively by conditioning on the last frame of the preceding chunk. 
%Moreover, we observed that  all image-to-video methods that we evaluated in our experiments (see \secrefcustom{sec:exp_comparison_with_baselines}) lead eventually to video stagnation,  when applied autoregressively by conditioning on the last frame of the previous chunk. 

To overcome the weaknesses of current works, we propose \textbf{\textit{StreamingT2V}}, an autoregressive text-to-video method equipped with long/short-term memory blocks that  generates long videos without temporal inconsistencies.

To this end, we propose the \textbf{\textit{Conditional Attention Module (CAM)}} which, due to its attentional nature, effectively borrows the content information from the previous frames to generate new ones, while not restricting their motion by the previous structures/shapes.
Thanks to CAM, our results are smooth and with artifact-free video chunk transitions. 

Current methods not only exhibit temporal inconsistencies and video stagnation, but also experience alterations in object appearance/characteristics (see \eg SEINE \cite{chen2023seine} in 
%% EXTERNAL REFERENCE ATTENTION
% \figrefcustom{appendix:fig:testset-compare-2})
Fig. \textcolor{linkblue}{A.4})
and a decline in video quality over time (see \eg SVD \cite{SVD} in \figrefcustom{fig:comparison-to-sota}).
This occurs as  only the last frame(s) of the preceding chunk are considered, thus overlooking long-term dependencies in the autoregressive process. 
%
%The reason is that, due to conditioning only on the last frame(s) of the previous chunk, they overlook the long-term dependencies of the autoregressive process. 
To address this issue we design an \textbf{\textit{Appearance Preservation Module (APM)}} that extracts object and global scene details from an initial image, to condition the video generation with that information, ensuring consistency in object and scene features  throughout the autoregressive process.   

To further enhance the quality and resolution of our long video generation, we adapt a video enhancement model for autoregressive generation.
To this end, we apply the SDEdit \cite{meng2021sdedit} approach on a high-resolution text-to-video model and  enhance consecutive 24-frame chunks (overlapping with 8 frames) of our video.
To make the chunk enhancement transitions smooth, we design a \textbf{\textit{randomized blending}} approach for seamless merging of overlapping chunks.

Experiments show that StreamingT2V generates long and temporal consistent videos from text without video stagnation. 
To summarize, our contributions are three-fold:
\begin{itemize}
    \item We introduce \textbf{\textit{StreamingT2V}}, an autoregressive approach for seamless synthesis of extended video content 
    using short and long-term dependencies. 
    \item Our \textbf{\textit{Conditional Attention Module (CAM)}} and \textbf{\textit{Appearance Preservation Module (APM)}} ensure the natural continuity of the global scene and object characteristics of generated videos.
    \item We seamlessly enhance generated long videos by introducing our \textbf{\textit{randomized blending approach}} of consecutive overlapping chunks.
\end{itemize}